\documentclass[11pt,a4paper]{article}

\usepackage{amsmath,amssymb,amsthm,mathrsfs}
\usepackage[margin=1in]{geometry}
\usepackage{hyperref}

%% Theorem environments
\theoremstyle{plain}
\newtheorem{theorem}{Theorem}[section]
\newtheorem{lemma}[theorem]{Lemma}
\newtheorem{proposition}[theorem]{Proposition}
\newtheorem{corollary}[theorem]{Corollary}

\theoremstyle{definition}
\newtheorem{definition}[theorem]{Definition}
\newtheorem{remark}[theorem]{Remark}

%% Notation
\newcommand{\T}{\mathbb{T}}
\newcommand{\R}{\mathbb{R}}
\newcommand{\cD}{\mathcal{D}}
\newcommand{\cH}{\mathcal{H}}
\newcommand{\cV}{\mathcal{V}}
\newcommand{\E}{\mathbb{E}}
\newcommand{\ip}[2]{\langle #1,\,#2\rangle}
\newcommand{\norm}[1]{\|#1\|}
% Wick product notation
\newcommand{\wick}[1]{\mathopen{:}\!#1\!\mathclose{:}}

\begin{document}

\title{Equivalence of the $\Phi^4_3$ Measure under Smooth Shifts}
\author{}
\date{}
\maketitle

\begin{abstract}
We prove that the $\Phi^4_3$ probability measure $\mu$ on the space of distributions
$\cD'(\T^3)$ is equivalent (mutually absolutely continuous) to its pushforward $T_\psi^*\mu$
under the shift $T_\psi(u)=u+\psi$, for every smooth function $\psi\in C^\infty(\T^3)$.
The proof combines the Cameron--Martin theorem for the Gaussian free field with an explicit
Wick-polynomial expansion of the shifted interaction.  The Radon--Nikodym derivative is
computed explicitly and shown to be $\mu$-a.s.\ strictly positive, establishing equivalence.
\end{abstract}

%% ================================================================
\section{Notation and Setup}
%% ================================================================

\subsection{Three-Dimensional Torus and Gaussian Free Field}

Let $\T^3=(\R/\Z)^3$ be the three-dimensional unit torus with normalised Lebesgue measure $dx$.
Denote by $\cD'(\T^3)$ the space of real-valued Schwartz distributions on $\T^3$, equipped with
its Borel $\sigma$-algebra for the weak-$*$ topology.

Fix a mass parameter $m>0$.  The \emph{Gaussian free field} is the centered Gaussian Borel
probability measure $\mu_0$ on $\cD'(\T^3)$ with covariance operator
\begin{equation}\label{eq:covariance}
  C := (m^2 - \Delta)^{-1} : L^2(\T^3) \longrightarrow L^2(\T^3),
\end{equation}
characterized by
\begin{equation}\label{eq:char}
  \int_{\cD'} e^{i\ip{f}{u}}\,d\mu_0(u)
  = \exp\!\Bigl(-\tfrac{1}{2}\ip{f}{Cf}_{L^2(\T^3)}\Bigr),
  \qquad f\in C^\infty(\T^3),
\end{equation}
where $\ip{f}{u}$ denotes the standard distribution pairing.
Existence follows from Minlos's theorem; the measure is supported on $H^s(\T^3)$
for every $s < -\tfrac{1}{2}$.

\subsection{Wick Powers}

In three dimensions the diagonal $C(x,x)=\sum_{k\in\mathbb{Z}^3}(m^2+4\pi^2|k|^2)^{-1}$ is
divergent, so the pointwise power $u(x)^n$ requires renormalization.

\begin{definition}[Wick powers]\label{def:wick}
For $n\ge 0$ and $f\in L^2(\T^3)$, define the \emph{smeared $n$-th Wick power}
\begin{equation}\label{eq:wick-def}
  \wick{u^n}_{\mu_0}(f) := I_n(f^{\otimes n}) \in L^2(\mu_0),
\end{equation}
where $I_n$ denotes the $n$-th Wiener--It\^o multiple stochastic integral with respect to $\mu_0$.
We write $\int_{\T^3}f(x)\wick{u(x)^n}_{\mu_0}\,dx := I_n(f^{\otimes n})$.

The Wick power $\wick{u(x)^n}_{\mu_0}$ is alternatively characterised by the regularised limit:
for any ultraviolet cutoff $u_\varepsilon = \rho_\varepsilon * u$ (with $\rho_\varepsilon$
a smooth approximate identity),
\[
  \wick{u_\varepsilon(x)^n}_{\mu_0} := H_n\!\bigl(u_\varepsilon(x);\,\sigma_\varepsilon^2(x)\bigr),
  \quad \sigma_\varepsilon^2(x) := \E_{\mu_0}[u_\varepsilon(x)^2],
\]
where $H_n(\,\cdot\,;\sigma^2)$ is the $n$-th probabilist Hermite polynomial with variance
parameter $\sigma^2$.  The limit $\varepsilon\to 0$ of $\int_{\T^3}f(x)\wick{u_\varepsilon(x)^n}_{\mu_0}\,dx$
converges in $L^2(\mu_0)$ and is independent of the choice of regularisation.
\end{definition}

We use the shorthand $\int_{\T^3}\wick{u^n}_{\mu_0}\,dx := \wick{u^n}_{\mu_0}(\mathbf{1}_{\T^3})$
for the constant test function $\mathbf{1}_{\T^3}$.

\subsection{The $\Phi^4_3$ Measure}

\begin{definition}[$\Phi^4_3$ measure]\label{def:phi4}
The \emph{$\Phi^4_3$ measure} is the Borel probability measure
\begin{equation}\label{eq:phi4-def}
  d\mu(u) = \frac{1}{Z}\,e^{-\cV(u)}\,d\mu_0(u),
\end{equation}
where the renormalized interaction is
\begin{equation}\label{eq:action}
  \cV(u) := \int_{\T^3}\wick{u(x)^4}_{\mu_0}\,dx
            + \delta m^2\int_{\T^3}\wick{u(x)^2}_{\mu_0}\,dx,
\end{equation}
with $\delta m^2\in\R$ a renormalization constant chosen (via the construction) so that
\[
  Z := \int_{\cD'} e^{-\cV(u)}\,d\mu_0(u) \in (0,\infty).
\]
\end{definition}

\begin{remark}\label{rem:construction}
The rigorous construction of $\mu$ -- including the definition of $\delta m^2$, the
$L^2(\mu_0)$-convergence of regularised versions of $\cV$, and the bound $Z\in(0,\infty)$
-- is established in Glimm--Jaffe~\cite{GJ87} via cluster expansions and in
Barashkov--Gubinelli~\cite{BG20} via a variational method.  Both references prove the
stronger statement $e^{-\cV}\in\bigcap_{p\ge 1}L^p(\mu_0)$; in particular
$\norm{e^{-p\cV}}_{L^1(\mu_0)}<\infty$ for all $p\ge 1$.  Since $e^{-\cV}>0$ everywhere,
the density $e^{-\cV}/Z$ is strictly positive $\mu_0$-a.s., so $\mu\sim\mu_0$.
\end{remark}

\subsection{The Shift Map}

For $\psi\in C^\infty(\T^3)$, define the \emph{shift} $T_\psi:\cD'(\T^3)\to\cD'(\T^3)$ by
$T_\psi(u) := u + \psi$, embedding $\psi$ in $\cD'(\T^3)$ via the inclusion
$C^\infty(\T^3)\hookrightarrow\cD'(\T^3)$.  The \emph{pushforward} is
\[
  (T_\psi^*\mu)(A) := \mu\bigl(T_\psi^{-1}(A)\bigr) = \mu(A-\psi),
  \qquad A \subset \cD'(\T^3)\text{ Borel},
\]
i.e., $T_\psi^*\mu$ is the distribution of $u+\psi$ when $u\sim\mu$.

%% ================================================================
\section{Statement of the Main Result}
%% ================================================================

\begin{theorem}\label{thm:main}
For every smooth function $\psi\in C^\infty(\T^3)$, the measures $\mu$ and $T_\psi^*\mu$
on $\cD'(\T^3)$ are equivalent (mutually absolutely continuous): $\mu \sim T_\psi^*\mu$.

More precisely, $T_\psi^*\mu \ll \mu$ with Radon--Nikodym derivative
\begin{equation}\label{eq:RN}
  \frac{d(T_\psi^*\mu)}{d\mu}(u)
  = \exp\!\Bigl(\cV(u) - \cV(u-\psi) + \hat{\psi}(u) - \tfrac{1}{2}\norm{\psi}_{\cH}^2\Bigr),
\end{equation}
where
\begin{equation}\label{eq:hatpsi}
  \hat{\psi}(u) := \ip{(m^2-\Delta)\psi}{u}_{C^\infty\times\cD'},
  \qquad
  \norm{\psi}_{\cH}^2 := \ip{\psi}{(m^2-\Delta)\psi}_{L^2(\T^3)}.
\end{equation}
This derivative is $\mu$-a.s.\ strictly positive, which together with $T_\psi^*\mu\ll\mu$
yields $\mu \ll T_\psi^*\mu$, hence $\mu \sim T_\psi^*\mu$.
\end{theorem}

%% ================================================================
\section{Key Tools}
%% ================================================================

\subsection{Cameron--Martin Theorem for $\mu_0$}\label{sec:CM}

\begin{theorem}[Cameron--Martin; \cite{Bog98} Theorem~3.2.2]\label{thm:CM}
Let $\gamma$ be a centered Gaussian Borel probability measure on a separable Hilbert (or
locally convex) space $X$, with covariance operator $C_\gamma$.
Define the \emph{Cameron--Martin space}
$\cH(\gamma) := C_\gamma^{1/2}(X)$
with inner product $\ip{h}{k}_{\cH(\gamma)} = \ip{C_\gamma^{-1/2}h}{C_\gamma^{-1/2}k}_{X}$.

For every $h \in \cH(\gamma)$:
\begin{enumerate}
\item[\textup{(a)}] $T_h^*\gamma \sim \gamma$ with Radon--Nikodym derivative
\begin{equation}\label{eq:CM-RN}
  \frac{d(T_h^*\gamma)}{d\gamma}(x)
  = \exp\!\Bigl(\hat{h}(x) - \tfrac{1}{2}\norm{h}_{\cH(\gamma)}^2\Bigr),
\end{equation}
where $\hat{h}$ is the element of $L^2(\gamma)$ satisfying
$\E_\gamma[\hat{h}(x)\,\ip{f}{x}] = \ip{f}{h}$ for all $f\in X^*$.

\item[\textup{(b)}] For every measurable $g:X\to[0,\infty]$:
\begin{equation}\label{eq:CM-sub}
  \int_X g(x + h)\,d\gamma(x)
  = \int_X g(x)\,e^{\hat{h}(x) - \frac{1}{2}\norm{h}_\cH^2}\,d\gamma(x).
\end{equation}
\end{enumerate}
\end{theorem}

\begin{lemma}[Cameron--Martin space of $\mu_0$]\label{lem:CM-space}
The Cameron--Martin space of $\mu_0$ is
\[
  \cH(\mu_0) = H^1(\T^3)
  \quad\text{with}\quad
  \norm{\psi}_{\cH(\mu_0)}^2 = \ip{\psi}{(m^2-\Delta)\psi}_{L^2(\T^3)}.
\]
In particular, $C^\infty(\T^3)\subset\cH(\mu_0)$.  For $\psi\in C^\infty(\T^3)$,
the Riesz representative is
\[
  \hat{\psi}(u) = \ip{(m^2-\Delta)\psi}{\,u\,}_{C^\infty\times\cD'},
\]
which is $\mu_0$-a.s.\ well-defined and is a Gaussian random variable with
$\E_{\mu_0}[\hat{\psi}] = 0$ and $\E_{\mu_0}[\hat{\psi}^2] = \norm{\psi}_{\cH}^2$.
\end{lemma}

\begin{proof}
The covariance operator is $C_{\mu_0}=(m^2-\Delta)^{-1}$, so
$C_{\mu_0}^{1/2}=(m^2-\Delta)^{-1/2}$.  In the Fourier basis
$\{e^{2\pi i k\cdot x}\}_{k\in\mathbb{Z}^3}$ of $L^2(\T^3)$:
\[
  (m^2-\Delta)^{-1/2}\hat{f}(k) = (m^2+4\pi^2|k|^2)^{-1/2}\hat{f}(k).
\]
The image $C_{\mu_0}^{1/2}(L^2(\T^3)) = (m^2-\Delta)^{-1/2}L^2(\T^3)$ consists
precisely of functions $\phi$ satisfying
$\sum_{k\in\mathbb{Z}^3}(m^2+4\pi^2|k|^2)|\hat\phi(k)|^2 < \infty$,
which is exactly the Sobolev space $H^1(\T^3)$.
Hence $\cH(\mu_0)=H^1(\T^3)$.

The Cameron--Martin inner product is
$\norm{\psi}_{\cH}^2 = \norm{(m^2-\Delta)^{1/2}\psi}_{L^2}^2
= \ip{(m^2-\Delta)^{1/2}\psi}{(m^2-\Delta)^{1/2}\psi}_{L^2}
= \ip{\psi}{(m^2-\Delta)\psi}_{L^2}$.

Since $C^\infty(\T^3)\subset H^k(\T^3)$ for all $k\ge0$, in particular
$C^\infty(\T^3)\subset H^1(\T^3)=\cH(\mu_0)$.

For the Riesz representative: $\hat\psi = C_{\mu_0}^{-1}\psi = (m^2-\Delta)\psi$,
paired with $u$ via the distribution pairing.  Since
$(m^2-\Delta)\psi\in C^\infty(\T^3)$ whenever $\psi\in C^\infty(\T^3)$,
the pairing $\ip{(m^2-\Delta)\psi}{u}$ is well-defined for every $u\in\cD'(\T^3)$.
The Gaussian character and moment formulas follow from the general Cameron--Martin
theory.
\end{proof}

\subsection{Wick Product Shift Formula}

\begin{lemma}[Wick shift formula]\label{lem:wick-shift}
For $\psi\in C^\infty(\T^3)$, $f\in L^2(\T^3)$, and $n\in\{0,1,2,3,4\}$,
the following identity holds in $L^2(\mu_0)$:
\begin{equation}\label{eq:wick-shift}
  \int_{\T^3} f(x)\,\wick{(u(x)-\psi(x))^n}_{\mu_0}\,dx
  = \sum_{k=0}^{n}\binom{n}{k}(-1)^{n-k}
    \int_{\T^3} f(x)\,\psi(x)^{n-k}\,\wick{u(x)^k}_{\mu_0}\,dx.
\end{equation}
In particular, with $f\equiv\mathbf{1}$:
\begin{align}
  \int_{\T^3}\wick{(u-\psi)^4}_{\mu_0}\,dx
  &= \int_{\T^3}\wick{u^4}_{\mu_0}\,dx
    - 4\int_{\T^3}\psi\,\wick{u^3}_{\mu_0}\,dx
    + 6\int_{\T^3}\psi^2\,\wick{u^2}_{\mu_0}\,dx \nonumber\\
  &\quad - 4\int_{\T^3}\psi^3\,u\,dx
    + \int_{\T^3}\psi^4\,dx, \label{eq:wick4}\\
  \int_{\T^3}\wick{(u-\psi)^2}_{\mu_0}\,dx
  &= \int_{\T^3}\wick{u^2}_{\mu_0}\,dx
    - 2\int_{\T^3}\psi\,u\,dx
    + \int_{\T^3}\psi^2\,dx. \label{eq:wick2}
\end{align}
\end{lemma}

\begin{proof}
At the regularised level, $\wick{u_\varepsilon(x)^n}_{\mu_0}=H_n(u_\varepsilon(x);\sigma_\varepsilon^2(x))$
where $H_n(\,\cdot\,;\sigma^2)$ is the probabilist Hermite polynomial.  We establish
the algebraic identity
\begin{equation}\label{eq:hermite-add}
  H_n(a+t;\,\sigma^2) = \sum_{k=0}^{n}\binom{n}{k}t^{n-k}\,H_k(a;\,\sigma^2)
\end{equation}
by comparing coefficients of $\alpha^n$ in the generating-function equality
\[
  \underbrace{e^{\alpha(a+t)-\frac{\alpha^2\sigma^2}{2}}}_{=\sum_n\frac{\alpha^n}{n!}H_n(a+t;\sigma^2)}
  = e^{\alpha t}\cdot\underbrace{e^{\alpha a - \frac{\alpha^2\sigma^2}{2}}}_{=\sum_k\frac{\alpha^k}{k!}H_k(a;\sigma^2)}
  = \Bigl(\sum_{j\ge0}\frac{(\alpha t)^j}{j!}\Bigr)\Bigl(\sum_{k\ge0}\frac{\alpha^k}{k!}H_k(a;\sigma^2)\Bigr).
\]
Extracting the coefficient of $\alpha^n$ on the right gives
$\frac{1}{n!}H_n(a+t;\sigma^2) = \sum_{k=0}^{n}\frac{t^{n-k}}{(n-k)!}\frac{H_k(a;\sigma^2)}{k!}$,
which is \eqref{eq:hermite-add}.

Apply \eqref{eq:hermite-add} at each $x\in\T^3$ with $a=u_\varepsilon(x)$ and $t=-\psi(x)$
(a deterministic quantity at each $x$, so there are no stochastic contractions between $\psi(x)$
and $u_\varepsilon(x)$):
\[
  \wick{(u_\varepsilon(x)-\psi(x))^n}_{\mu_0}
  = H_n(u_\varepsilon(x)-\psi(x);\sigma_\varepsilon^2(x))
  = \sum_{k=0}^{n}\binom{n}{k}(-\psi(x))^{n-k}\,H_k(u_\varepsilon(x);\sigma_\varepsilon^2(x)).
\]
Multiplying by $f(x)$ and integrating over $\T^3$, then taking $\varepsilon\to0$ in $L^2(\mu_0)$
(the limit is justified since $f\psi^{n-k}\in L^2(\T^3)$ for smooth $\psi$ and $f\in L^2$),
we obtain \eqref{eq:wick-shift}.

The identities \eqref{eq:wick4}--\eqref{eq:wick2} follow from \eqref{eq:wick-shift}
with $f\equiv\mathbf{1}$ and $n=4,2$ respectively, using
$H_0(\,\cdot\,;\sigma^2)=1$ and $H_1(v;\sigma^2)=v$,
so $\wick{u^0}_{\mu_0}=1$ and $\wick{u^1}_{\mu_0}=u$.
\end{proof}

\subsection{Integrability of Smeared Wick Polynomials Under $\mu$}

\begin{lemma}[Wick integrability under $\mu$]\label{lem:integrability}
For every $k\in\{0,1,2,3\}$, every $f\in C^\infty(\T^3)$, and every $p\in[1,\infty)$:
\begin{equation}\label{eq:Lp-bound}
  W_k := \int_{\T^3} f(x)\,\wick{u(x)^k}_{\mu_0}\,dx
  \;\in\; L^p(\mu).
\end{equation}
\end{lemma}

\begin{proof}
By H\"{o}lder's inequality with exponents $2$ and $2$:
\begin{equation}\label{eq:holder}
  \norm{W_k}_{L^p(\mu)}^p
  = \frac{1}{Z}\int|W_k|^p e^{-\cV(u)}\,d\mu_0(u)
  \le \frac{1}{Z}\,\norm{W_k}_{L^{2p}(\mu_0)}^{p}
      \cdot\norm{e^{-\cV}}_{L^2(\mu_0)}.
\end{equation}

\emph{Bounding $\norm{W_k}_{L^{2p}(\mu_0)}$:}
The element $W_k = I_k(f^{\otimes k})$ lies in the $k$-th Wiener--It\^o chaos of $\mu_0$.
Nelson's hypercontractivity theorem~\cite{Nel73} states that for any $q\ge 2$ and any
element $F$ of the $k$-th chaos,
$\norm{F}_{L^q(\mu_0)} \le (q-1)^{k/2}\norm{F}_{L^2(\mu_0)}$.
Applying this with $q=2p$ gives
$\norm{W_k}_{L^{2p}(\mu_0)} \le (2p-1)^{k/2}\norm{W_k}_{L^2(\mu_0)} < \infty$,
since $f\in L^2(\T^3)$ implies $\norm{W_k}_{L^2(\mu_0)}
= \norm{f}_{L^2(\T^3)}^k \cdot k! \cdot \norm{C^{k/2}(\,\cdot\,)}^k < \infty$
(the $L^2$ norm of a Wiener chaos element is controlled by the $L^2$ norm of its kernel).

\emph{Bounding $\norm{e^{-\cV}}_{L^2(\mu_0)}$:}
The $\Phi^4_3$ construction (see Remark~\ref{rem:construction} and \cite{BG20}, Proposition~3.1)
establishes $e^{-\cV}\in\bigcap_{p\ge1}L^p(\mu_0)$.
In particular $\norm{e^{-\cV}}_{L^2(\mu_0)}<\infty$.

Combining both bounds in \eqref{eq:holder} gives $\norm{W_k}_{L^p(\mu)}<\infty$.
\end{proof}

%% ================================================================
\section{Proof of Theorem~\ref{thm:main}}
%% ================================================================

\begin{proof}[Proof of Theorem~\ref{thm:main}]

\medskip\noindent
\textbf{Step~1: Density of $T_\psi^*\mu$ with respect to $\mu_0$.}

For every Borel measurable $h:\cD'(\T^3)\to[0,\infty)$, the pushforward formula gives
\begin{align}
  \int h\,d(T_\psi^*\mu)
  &= \int h(T_\psi(v))\,d\mu(v)
   = \int h(v+\psi)\,\frac{e^{-\cV(v)}}{Z}\,d\mu_0(v). \label{eq:step1-start}
\end{align}
We convert the right-hand side to an integral against $\mu_0$ without the shift.
Define $F:\cD'\to[0,\infty)$ by
$F(u) := h(u)\,e^{-\cV(u-\psi)}/Z$,
so that $F(T_\psi(v)) = F(v+\psi) = h(v+\psi)\,e^{-\cV(v)}/Z$.
By the standard pushforward formula,
\[
  \int h(v+\psi)\,\frac{e^{-\cV(v)}}{Z}\,d\mu_0(v)
  = \int F(T_\psi(v))\,d\mu_0(v)
  = \int F(u)\,d(T_\psi^*\mu_0)(u).
\]
Applying Lemma~\ref{lem:CM-space} ($\psi\in\cH(\mu_0)$) and
Theorem~\ref{thm:CM}(a) ($d(T_\psi^*\mu_0) = e^{\hat\psi - \frac{1}{2}\norm\psi_\cH^2}\,d\mu_0$):
\begin{align}
  \int F(u)\,d(T_\psi^*\mu_0)(u)
  &= \int h(u)\,\frac{e^{-\cV(u-\psi)}}{Z}\,
     e^{\hat\psi(u)-\frac{1}{2}\norm\psi_\cH^2}\,d\mu_0(u). \label{eq:step1-end}
\end{align}
Since \eqref{eq:step1-start}--\eqref{eq:step1-end} hold for all measurable $h\ge0$,
we conclude $T_\psi^*\mu\ll\mu_0$ with density
\begin{equation}\label{eq:density-mu0}
  \frac{d(T_\psi^*\mu)}{d\mu_0}(u)
  = \frac{e^{-\cV(u-\psi)}}{Z}\,e^{\hat\psi(u)-\frac{1}{2}\norm\psi_\cH^2}.
\end{equation}

\medskip\noindent
\textbf{Step~2: Normalization.}

We verify that the right-hand side of \eqref{eq:density-mu0} integrates to~$1$ against $\mu_0$.
Apply the Cameron--Martin substitution formula \eqref{eq:CM-sub}
(Theorem~\ref{thm:CM}(b)) to $\gamma=\mu_0$, shift $\psi$, and
test function $g:\cD'\to[0,\infty)$ defined by
$g(x) := e^{-\cV(x-\psi)}/Z$.
Since $g(x+\psi) = e^{-\cV(x)}/Z$, formula \eqref{eq:CM-sub} gives:
\begin{equation}\label{eq:norm-check}
  \int_{\cD'}\frac{e^{-\cV(u-\psi)}}{Z}\,e^{\hat\psi(u)-\frac{1}{2}\norm\psi_\cH^2}\,d\mu_0(u)
  = \int_{\cD'}\frac{e^{-\cV(u)}}{Z}\,d\mu_0(u)
  = \frac{Z}{Z} = 1.
\end{equation}
This confirms that \eqref{eq:density-mu0} is a valid probability density.

\medskip\noindent
\textbf{Step~3: Radon--Nikodym derivative with respect to $\mu$.}

Since $\mu \sim \mu_0$ (Remark~\ref{rem:construction}), we apply the chain rule:
\begin{equation}\label{eq:chain-rule}
  \frac{d(T_\psi^*\mu)}{d\mu}(u)
  = \frac{d(T_\psi^*\mu)}{d\mu_0}(u)
    \cdot\frac{d\mu_0}{d\mu}(u).
\end{equation}
The inverse density is $d\mu_0/d\mu = (d\mu/d\mu_0)^{-1} = Z\,e^{\cV(u)}$,
valid $\mu$-a.s.\ since $e^{-\cV(u)}>0$ everywhere.
Substituting into \eqref{eq:chain-rule}:
\begin{align}
  \frac{d(T_\psi^*\mu)}{d\mu}(u)
  &= \frac{e^{-\cV(u-\psi)}}{Z}\,e^{\hat\psi(u)-\frac{1}{2}\norm\psi_\cH^2}
     \cdot Z\,e^{\cV(u)} \notag\\
  &= \exp\!\Bigl(\cV(u)-\cV(u-\psi)+\hat\psi(u)-\tfrac{1}{2}\norm\psi_\cH^2\Bigr).
  \label{eq:RN-explicit}
\end{align}
This is precisely the formula stated in \eqref{eq:RN}.

\medskip\noindent
\textbf{Step~4: Expansion of $\cV(u)-\cV(u-\psi)$.}

Using Lemma~\ref{lem:wick-shift} (equations \eqref{eq:wick4}--\eqref{eq:wick2}):
\begin{align}
  \cV(u)-\cV(u-\psi)
  &= \Bigl[\int_{\T^3}\wick{u^4}_{\mu_0}\,dx
           - \int_{\T^3}\wick{(u-\psi)^4}_{\mu_0}\,dx\Bigr] \notag\\
  &\quad + \delta m^2\Bigl[\int_{\T^3}\wick{u^2}_{\mu_0}\,dx
           - \int_{\T^3}\wick{(u-\psi)^2}_{\mu_0}\,dx\Bigr] \notag\\[4pt]
  &= 4\!\int_{\T^3}\!\psi\,\wick{u^3}_{\mu_0}\,dx
     - 6\!\int_{\T^3}\!\psi^2\wick{u^2}_{\mu_0}\,dx
     + 4\!\int_{\T^3}\!\psi^3\,u\,dx
     - \!\int_{\T^3}\!\psi^4\,dx \notag\\
  &\quad + 2\delta m^2\!\int_{\T^3}\!\psi\,u\,dx
     - \delta m^2\!\int_{\T^3}\!\psi^2\,dx.
  \label{eq:V-diff}
\end{align}

We verify that each term in \eqref{eq:V-diff} lies in $L^p(\mu)$ for every $p\ge1$:
\begin{itemize}
\item[(a)]
  $\int_{\T^3}\psi^4\,dx$ and $\delta m^2\int_{\T^3}\psi^2\,dx$ are finite real constants
  (since $\psi\in C^\infty\subset L^4\cap L^2$).
\item[(b)]
  $\int_{\T^3}\psi\,u\,dx = \ip{\psi}{u}_{C^\infty\times\cD'}$ and
  $\int_{\T^3}\psi^3 u\,dx = \ip{\psi^3}{u}$:
  both are in $L^p(\mu)$ for all $p$ by Lemma~\ref{lem:integrability}
  with $k=1$, $f=\psi\in C^\infty$ and $f=\psi^3\in C^\infty$ respectively.
\item[(c)]
  $\int_{\T^3}\psi^2\wick{u^2}_{\mu_0}\,dx$: in $L^p(\mu)$ for all $p$
  by Lemma~\ref{lem:integrability} with $k=2$ and $f=\psi^2\in C^\infty$.
\item[(d)]
  $\int_{\T^3}\psi\,\wick{u^3}_{\mu_0}\,dx$: in $L^p(\mu)$ for all $p$
  by Lemma~\ref{lem:integrability} with $k=3$ and $f=\psi\in C^\infty$.
\end{itemize}
Hence $\cV(u)-\cV(u-\psi)\in L^p(\mu)$ for every $p\ge1$,
and in particular is $\mu$-a.s.\ finite.

Similarly, $\hat\psi(u)=\ip{(m^2-\Delta)\psi}{u}\in L^p(\mu)$ for all $p$
by Lemma~\ref{lem:integrability} with $k=1$ and $f=(m^2-\Delta)\psi\in C^\infty$,
and $\norm\psi_\cH^2<\infty$ since $\psi\in H^1(\T^3)$.

\medskip\noindent
\textbf{Step~5: Strict positivity and equivalence.}

The exponent in \eqref{eq:RN-explicit} is $\mu$-a.s.\ finite (Step~4), so the
Radon--Nikodym derivative $d(T_\psi^*\mu)/d\mu = e^{(\cdots)}$ is $\mu$-a.s.\ strictly
positive.  This establishes $T_\psi^*\mu \ll \mu$ (Step~3 and Step~2).

For the reverse inclusion: let $A\subset\cD'$ be any Borel set with $(T_\psi^*\mu)(A)=0$.
Then
\[
  0 = (T_\psi^*\mu)(A)
    = \int_A \frac{d(T_\psi^*\mu)}{d\mu}(u)\,d\mu(u)
    = \int_A e^{\cV(u)-\cV(u-\psi)+\hat\psi(u)-\frac{1}{2}\norm\psi_\cH^2}\,d\mu(u).
\]
Since the integrand is $\mu$-a.s.\ strictly positive and the integral over $A$ vanishes,
we conclude $\mu(A)=0$.  Hence $\mu \ll T_\psi^*\mu$.

Therefore $\mu \sim T_\psi^*\mu$, completing the proof.
\end{proof}

%% ================================================================
\section{Verification of Hypotheses}
%% ================================================================

We systematically verify every precondition invoked in the proof.

\begin{enumerate}
\item \textbf{Cameron--Martin space inclusion} (Lemma~\ref{lem:CM-space}).
  We need $\psi\in\cH(\mu_0)=H^1(\T^3)$.
  Since $\psi\in C^\infty(\T^3)$, we have $\psi\in H^k(\T^3)$ for all $k\ge0$,
  in particular $\psi\in H^1(\T^3)$.  \textit{Verified.}

\item \textbf{Cameron--Martin Theorem~\ref{thm:CM} applies.}
  Requires $\psi\in\cH(\mu_0)$.  Verified in (1).  \textit{Verified.}

\item \textbf{Riesz representative $\hat\psi$ is well-defined.}
  Requires $(m^2-\Delta)\psi$ to lie in $C^\infty(\T^3)$ (to pair with $u\in\cD'$).
  Since $\psi\in C^\infty(\T^3)$ and $m^2-\Delta$ is a differential operator with smooth
  coefficients on $\T^3$, $(m^2-\Delta)\psi\in C^\infty(\T^3)$.  \textit{Verified.}

\item \textbf{Hermite addition formula \eqref{eq:hermite-add}.}
  This is a purely algebraic identity of polynomials, valid for all $a,t\in\R$ and $\sigma^2>0$.
  Proved in Lemma~\ref{lem:wick-shift}.  \textit{Verified.}

\item \textbf{Wick shift formula passes through the $\varepsilon\to0$ limit.}
  At the regularised level the identity holds exactly (as an algebraic identity in $u_\varepsilon(x)$
  and the deterministic $\psi(x)$).  The limits $\int f\wick{u_\varepsilon^k}_{\mu_0}\,dx \to \int f\wick{u^k}_{\mu_0}\,dx$
  hold in $L^2(\mu_0)$ (by definition of the Wick powers) and $f\psi^{n-k}\in L^2(\T^3)$
  whenever $f\in L^2$ and $\psi\in C^\infty$.  \textit{Verified.}

\item \textbf{Hypercontractivity} (Lemma~\ref{lem:integrability}).
  Nelson's theorem~\cite{Nel73} requires $W_k$ to belong to the $k$-th Wiener--It\^o chaos
  with $k<\infty$.  Here $k\le 3$ and $f\in C^\infty\subset L^2$, so $W_k=I_k(f^{\otimes k})$
  is indeed a chaos-$k$ element.  \textit{Verified.}

\item \textbf{Exponential integrability} (Lemma~\ref{lem:integrability}).
  Requires $e^{-\cV}\in L^2(\mu_0)$.  The construction proves
  $e^{-\cV}\in\bigcap_{p\ge1}L^p(\mu_0)$ \cite{BG20}.  \textit{Verified.}

\item \textbf{Inverse density $d\mu_0/d\mu = Ze^{\cV}$.}
  Valid $\mu$-a.s.\ because $e^{-\cV(u)}>0$ for all $u$ (exponential is strictly positive),
  so $d\mu/d\mu_0 = e^{-\cV}/Z>0$ $\mu_0$-a.s., and hence $d\mu_0/d\mu = Ze^{\cV}$ $\mu$-a.s.
  \textit{Verified.}

\item \textbf{Partition function $Z\in(0,\infty)$.}
  Part of the $\Phi^4_3$ construction; see \cite{GJ87,BG20}.  \textit{Verified.}
\end{enumerate}

%% ================================================================
\section{Boundary and Degenerate Cases}
%% ================================================================

\begin{enumerate}
\item \textbf{$\psi\equiv 0$ (excluded by hypothesis).}
  The problem specifies $\psi\not\equiv0$.  Nevertheless our proof is valid for $\psi=0$:
  the Radon--Nikodym derivative \eqref{eq:RN} equals~$1$, and $T_0^*\mu=\mu$ trivially.
  The non-triviality of the equivalence lies in $\psi\not\equiv0$, which our proof
  handles uniformly.

\item \textbf{$\psi$ with zero spatial mean $\int_{\T^3}\psi\,dx=0$.}
  No special treatment needed.  The Cameron--Martin theorem applies to all
  $\psi\in H^1(\T^3)$, with or without the zero-mean condition.

\item \textbf{$\psi$ with large $H^1$ norm.}
  The Cameron--Martin R-N derivative satisfies
  $\norm{e^{\hat\psi-\norm\psi_\cH^2/2}}_{L^2(\mu_0)} = e^{\norm\psi_\cH^2/2}$
  (since $\hat\psi\sim N(0,\norm\psi_\cH^2)$ under $\mu_0$,
  so $\E_{\mu_0}[e^{2\hat\psi-\norm\psi^2_\cH}]=e^{2\norm\psi^2_\cH - \norm\psi^2_\cH}=e^{\norm\psi^2_\cH}$,
  giving $\norm{e^{\hat\psi-\norm\psi^2_\cH/2}}_{L^2(\mu_0)}=e^{\norm\psi^2_\cH/2}$).
  This is finite for any fixed $\psi\in H^1(\T^3)$, regardless of $\norm\psi_\cH$.

\item \textbf{Sign of individual Wick monomials.}
  The Wick powers $\wick{u^4}_{\mu_0}$ and $\wick{u^2}_{\mu_0}$ are not sign-definite
  (Hermite polynomials of degree $\ge2$ are not positive).  This does not create any
  difficulty: the interaction $\cV$ is well-defined in $L^2(\mu_0)$ by the $\Phi^4_3$
  construction, and the difference $\cV(u)-\cV(u-\psi)$ given by \eqref{eq:V-diff}
  is a finite sum of $L^p(\mu)$ random variables (Step~4 of the proof).
\end{enumerate}

%% ================================================================
\section{Conclusion}
%% ================================================================

The answer to the problem is \textbf{yes}: the $\Phi^4_3$ measure $\mu$ and its
pushforward $T_\psi^*\mu$ under any smooth shift are equivalent probability measures
on $\cD'(\T^3)$.

The mechanism is transparent: a smooth function $\psi\in C^\infty(\T^3)$ belongs to
the Cameron--Martin space $H^1(\T^3)$ of the Gaussian reference measure $\mu_0$;
the $\Phi^4_3$ measure is an absolutely continuous perturbation of $\mu_0$; and when
the Gibbs weight is shifted by $\psi$, the change in interaction is a degree-$\le3$
polynomial in smeared Wick powers with smooth coefficients (equation~\eqref{eq:V-diff}),
which is $\mu$-integrable by hypercontractivity.  The resulting Radon--Nikodym
derivative \eqref{eq:RN} is $\mu$-a.s.\ strictly positive, giving equivalence in
both directions.

\begin{thebibliography}{10}

\bibitem{BG20}
N.~Barashkov and M.~Gubinelli,
\emph{A variational method for $\Phi^4_3$},
Duke Math.\ J.\ \textbf{169} (2020), no.~17, 3339--3415.

\bibitem{Bog98}
V.~I.~Bogachev,
\emph{Gaussian Measures},
Mathematical Surveys and Monographs, vol.~62,
American Mathematical Society, Providence, RI, 1998.

\bibitem{GJ87}
J.~Glimm and A.~Jaffe,
\emph{Quantum Physics: A Functional Integral Point of View},
2nd ed., Springer-Verlag, New York, 1987.

\bibitem{Nel73}
E.~Nelson,
\emph{The free Markov field},
J.\ Funct.\ Anal.\ \textbf{12} (1973), 211--227.

\end{thebibliography}

\end{document}
